%% 第五章 载人登月逃逸任务风险性评估

\chapter{载人登月逃逸任务风险性评估}
\chaptermark{载人登月逃逸任务风险性评估}

\section{评估指标体系}

基于第4章蒙特卡洛仿真的统计结果，本章建立针对发射逃逸阶段的风险评估体系~\cite{刁伟鹤2025,Antonsen2023}。评估流程分三步进行：首先依据拉丁超立方样本统计各关键指标的超限概率，然后对照三级风险等级判定单指标风险，最后结合正交试验的敏感性分析识别主要风险源，综合判断两种工况的整体风险水平。评估指标定义与阈值如表~\ref{tab:risk_threshold}~所示。

\begin{table}[htbp]
\caption{风险评估指标与阈值}
\label{tab:risk_threshold}
\begin{tabularx}{\textwidth}{CCCCC}
\toprule
指标 & 符号 & 零高度阈值 & 最大动压阈值 & 物理含义 \\
\midrule
水平脱靶量 & $D_h$ & $50\ \mathrm m$ & $1000\ \mathrm m$ & 末端落点偏离靶点的水平距离，超过则判定为脱靶 \\
三维脱靶量 & $D_{3D}$ & $100\ \mathrm m$ & $1500\ \mathrm m$ & 空间总偏差，辅助评估搜救区域覆盖 \\
最大角速度 & $\omega_{max}$ & $1.0\ \mathrm{rad/s}$ & $1.0\ \mathrm{rad/s}$ & ACM姿控系统角速度响应上限 \\
最高点高度 & $h_{max}$ & $\ge 1500\ \mathrm m$ & $\ge 1500\ \mathrm m$ & 保障降落伞展开序列的最低安全高度 \\
峰值动压 & $q_{max}$ & $\le 50\ \mathrm{kPa}$ & $\le 50\ \mathrm{kPa}$ & 飞行器结构及分离机构承载上限 \\
末端速度 & $V_f$ & 无独立阈值 & 无独立阈值 & 辅助参考，不开伞段末速不具备约束意义 \\
\bottomrule
\end{tabularx}
\end{table}

命中半径$50\ \mathrm{m}$（零高度）和$1000\ \mathrm{m}$（最大动压）直接采用第4章基于500组LHS样本$95\%$分位数确定的命中判据。角速度阈值$1.0\ \mathrm{rad/s}$由估算得到。最高点高度阈值$1500\ \mathrm{m}$参照猎户座PA-1试验经验：该试验远地点约$1800\ \mathrm{m}$，成功完成降落伞展开与着陆，因此$1500\ \mathrm{m}$可作为最低安全高度的参考下限。峰值动压阈值$50\ \mathrm{kPa}$也由估算得到。

风险等级按超限概率$P_{exceed}$划分为三级：

\begin{itemize}
    \item 低风险（$P_{exceed}\le5\%$）：超限概率低于$p95$分位数，系统在绝大多数扰动组合下满足指标要求；
    \item 中风险（$5\%< P_{exceed}\le15\%$）：超限概率在可接受范围内，需关注主要敏感性因素；
    \item 高风险（$P_{exceed}>15\%$）：超限概率偏高，系统对扰动较为敏感，需要重点分析。
\end{itemize}

\section{零高度逃逸风险评估}

\subsection{命中风险评估}

零高度逃逸500组LHS样本的命中统计结果如表~\ref{tab:pad_risk_hit}~所示。

\begin{table}[htbp]
\caption{零高度逃逸命中风险统计}
\label{tab:pad_risk_hit}
\begin{tabularx}{\textwidth}{>{\hsize=1.3\hsize}C >{\hsize=0.7\hsize}C}
\toprule
指标 & 数值 \\
\midrule
命中率 & $96.20\%$ \\
超限概率 $P_{exceed}$ & $3.80\%$（低风险） \\
$D_h$ 均值 & $25.473\ \mathrm m$ \\
$D_h$ 标准差 & $11.316\ \mathrm m$ \\
$D_h$ $p95$ 分位数 & $47.331\ \mathrm m$ \\
$D_h$ 最大值 & $69.820\ \mathrm m$ \\
\bottomrule
\end{tabularx}
\end{table}

命中风险等级为低风险。$96.20\%$ 的样本落在 $50\ \mathrm m$ 命中半径以内，$3.80\%$ 的超限比例低于 $5\%$ 的低风险阈值。最恶劣 $5\%$ 样本的脱靶量分布在 $47\sim70\ \mathrm m$ 区间，整体偏差幅度可控。

三维脱靶量 $D_{3D}$ 的统计结果为均值 $47.7\ \mathrm m$、最大值 $161.5\ \mathrm m$。以 $100\ \mathrm m$ 为阈值时超限19组，超限概率与水平脱靶量一致（$3.80\%$），表明零高度工况下的高度方向偏差不会引入额外的尾部风险，空间散布主要由水平散布主导。

\subsection{姿态风险评估}

最大角速度统计如表~\ref{tab:pad_risk_omega}~所示。

\begin{table}[htbp]
\caption{零高度逃逸姿态风险统计}
\label{tab:pad_risk_omega}
\begin{tabularx}{\textwidth}{>{\hsize=1.3\hsize}C >{\hsize=0.7\hsize}C}
\toprule
指标 & 数值 \\
\midrule
$\omega_{max}$ 均值 & $0.8674\ \mathrm{rad/s}$ \\
$\omega_{max}$ 最大值 & $0.8700\ \mathrm{rad/s}$ \\
变化幅度（最大值$-$均值） & $0.0025\ \mathrm{rad/s}$ \\
\bottomrule
\end{tabularx}
\end{table}

姿态风险等级为低风险。500组样本间的角速度变化幅度仅为 $0.0025\ \mathrm{rad/s}$，表明标称PD控制增益下ACM姿控系统对各类扰动均有充分的控制裕度，扰动不引起姿态失稳。标称最大角速度 $0.867\ \mathrm{rad/s}$ 主要由压缩制导时间线后的姿态重定向机动决定，与扰动无关。此外，所有样本的角速度均远低于 $1.0\ \mathrm{rad/s}$ 的工程阈值，角速度超限概率 $P_{exceed}=0\%$。

峰值动压 $q_{max}$ 的平均值为 $20.3\ \mathrm{kPa}$，最大值 $23.9\ \mathrm{kPa}$，远低于 $50\ \mathrm{kPa}$ 的结构承载阈值，超限概率 $P_{exceed}=0\%$。最高点高度 $h_{max}$ 的平均值为 $1782.5\ \mathrm m$，最小值 $1638.6\ \mathrm m$，全部500组样本均高于 $1500\ \mathrm m$ 的最低安全高度，超限概率 $P_{exceed}=0\%$，降落伞展开高度裕度充足。

\subsection{敏感性风险分析}

正交试验的极差分析揭示了各干扰因素对命中精度的贡献程度。零高度逃逸的敏感性排序如表~\ref{tab:pad_sens}~所示，风扰（特别是风向 $\chi_w$ 和水平风速 $u_0$）是影响脱靶量的第一敏感因素，极差 $R_{Dh}=13.00\ \mathrm m$；初始状态偏差紧随其后（$8.08\ \mathrm m$），质量偏差和推力幅值偏差居中（$5\sim6\ \mathrm m$）；发动机安装偏差的影响最小，极差仅 $1.11\ \mathrm m$。

\begin{table}[htbp]
\caption{零高度逃逸敏感性排序}
\label{tab:pad_sens}
\begin{tabularx}{\textwidth}{CCC}
\toprule
排名 & 因素 & 极差 $R_{Dh}$ (m) \\
\midrule
1 & 风扰-$\chi_w$（风向） & 13.00 \\
2 & 风扰-$u_0$（水平风速） & 8.10 \\
3 & 初始状态偏差 & 8.08 \\
4 & 质量偏差 & 6.01 \\
5 & 推力幅值偏差 & 5.47 \\
6 & 气动和大气模型偏差 & 4.77 \\
7 & 风扰-$w_0$（垂直风速） & 3.06 \\
8 & 发动机安装偏差 & 1.11 \\
\bottomrule
\end{tabularx}
\end{table}

从物理机制分析，零高度逃逸初始阶段速度低、气动阻尼弱，姿态控制主要依赖ACM喷气直接维持稳定。风场扰动作用于飞行器气动外形所产生的附加力和力矩无法被ACM完全抑制，因此风的强度和方向成为最主要的风险来源。

\section{最大动压逃逸风险评估}

\subsection{命中风险评估}

最大动压逃逸500组LHS样本的命中统计结果如表~\ref{tab:mq_risk_hit}~所示。

\begin{table}[htbp]
\caption{最大动压逃逸命中风险统计}
\label{tab:mq_risk_hit}
\begin{tabularx}{\textwidth}{>{\hsize=1.3\hsize}C >{\hsize=0.7\hsize}C}
\toprule
指标 & 数值 \\
\midrule
命中率 & $95.20\%$ \\
超限概率 $P_{exceed}$ & $4.80\%$（低风险） \\
$D_h$ 均值 & $456.333\ \mathrm m$ \\
$D_h$ 标准差 & $242.785\ \mathrm m$ \\
$D_h$ $p95$ 分位数 & $977.322\ \mathrm m$ \\
$D_h$ 最大值 & $1414.684\ \mathrm m$ \\
\bottomrule
\end{tabularx}
\end{table}

命中风险等级为低风险。$95.20\%$ 的样本落在 $1000\ \mathrm m$ 命中半径以内，超限比例 $4.80\%$。但与零高度逃逸相比，散布量级显著增大——$D_h$ 标准差达 $242.785\ \mathrm m$，最恶劣组合下的脱靶量可达 $1414.684\ \mathrm m$，尾部风险远高于零高度工况。

三维脱靶量 $D_{3D}$ 的统计结果为均值 $468.4\ \mathrm m$、最大值 $1444.2\ \mathrm m$，$p95$ 分位 $413.2\ \mathrm m$。以 $1500\ \mathrm m$ 为阈值时超限0组，超限概率 $P_{exceed}=0\%$。与零高度逃逸不同，最大动压工况的 $D_{3D}$ 最大值主要由水平分量贡献，高度方向的偏差占比相对较小。

\subsection{姿态风险评估}

最大角速度统计如表~\ref{tab:mq_risk_omega}~所示。

\begin{table}[htbp]
\caption{最大动压逃逸姿态风险统计}
\label{tab:mq_risk_omega}
\begin{tabularx}{\textwidth}{>{\hsize=1.3\hsize}C >{\hsize=0.7\hsize}C}
\toprule
指标 & 数值 \\
\midrule
$\omega_{max}$ 均值 & $0.2648\ \mathrm{rad/s}$ \\
$\omega_{max}$ 最大值 & $0.2652\ \mathrm{rad/s}$ \\
变化幅度（最大值$-$均值） & $0.0004\ \mathrm{rad/s}$ \\
\bottomrule
\end{tabularx}
\end{table}

姿态风险等级为低风险。最大角速度均值仅为零高度逃逸的三分之一（$0.2648\ \mathrm{rad/s}$ 对比 $0.8674\ \mathrm{rad/s}$），全样本变化幅度仅 $0.0004\ \mathrm{rad/s}$（约 $0.15\%$），角速度超限概率 $P_{exceed}=0\%$。物理上，高动压条件下的气动恢复力矩为姿态控制系统提供了充足的自然阻尼，因此最大动压逃逸的姿态稳定性优于零高度逃逸。

峰值动压 $q_{max}$ 的平均值为 $45.7\ \mathrm{kPa}$，最大值 $49.8\ \mathrm{kPa}$，$p95$ 分位 $45.7\ \mathrm{kPa}$。全部500组样本均低于 $50\ \mathrm{kPa}$ 的结构承载阈值，超限概率 $P_{exceed}=0\%$，但最大值已十分接近阈值（$49.8\ \mathrm{kPa}$），表明在极端扰动组合下峰值动压已接近结构设计的边界，属于值得关注的尾部风险。最高点高度 $h_{max}$ 的平均值为 $16450.0\ \mathrm m$，最小值 $16165.2\ \mathrm m$，全部样本均远高于 $1500\ \mathrm m$ 的最低安全高度，超限概率 $P_{exceed}=0\%$。

\subsection{敏感性风险分析}

最大动压逃逸的敏感性排序如表~\ref{tab:mq_sens}~所示。与零高度逃逸截然不同，气动和大气模型偏差是最核心的风险因素，极差 $R_{Dh}=276.30\ \mathrm m$，远高于其他因素。垂直风 $w_0$（$195.59\ \mathrm m$）和风向 $\chi_w$（$111.07\ \mathrm m$）也是重要影响因素。

\begin{table}[htbp]
\caption{最大动压逃逸敏感性排序}
\label{tab:mq_sens}
\begin{tabularx}{\textwidth}{CCC}
\toprule
排名 & 因素 & 极差 $R_{Dh}$ (m) \\
\midrule
1 & 气动和大气模型偏差 & 276.30 \\
2 & 风扰-$w_0$（垂直风速） & 195.59 \\
3 & 风扰-$\chi_w$（风向） & 111.07 \\
4 & 质量偏差 & 86.47 \\
5 & 发动机安装偏差 & 35.52 \\
6 & 初始状态偏差 & 35.32 \\
7 & 风扰-$u_0$（水平风速） & 27.95 \\
8 & 推力幅值偏差 & 14.38 \\
\bottomrule
\end{tabularx}
\end{table}

物理机制上，最大动压逃逸初始动压高（约30 kPa），气动参数的小幅偏差通过 $q=\frac12\rho V^2S$ 放大为显著的力和力矩偏差，且长达80 s的飞行时间使该偏差持续积累。因此气动系数和大气密度的工程不确定性成为制约弹道精度的首要风险。

\section{综合评估结论}

两种工况的完整风险评估结果对比如表~\ref{tab:risk_compare}~所示，超限概率汇总如表~\ref{tab:exceed_summary}~所示。

\begin{table}[htbp]
\caption{两种工况风险评估对比}
\label{tab:risk_compare}
\begin{tabularx}{\textwidth}{CCC}
\toprule
评估项 & 零高度逃逸 & 最大动压逃逸 \\
\midrule
命中风险 & 低风险（$3.80\%$） & 低风险（$4.80\%$） \\
姿态风险 & 低风险 & 低风险 \\
主要风险因素 & 风扰（风向/风速） & 气动和大气模型偏差 \\
次要风险因素 & 初始状态偏差、质量偏差 & 风扰（垂直风/风向） \\
尾部风险特征 & 紧凑（$\max\ D_h\approx 70\ \mathrm m$） & 分散（$\max\ D_h\approx 1415\ \mathrm m$） \\
\bottomrule
\end{tabularx}
\end{table}

\begin{table}[htbp]
\caption{超限概率汇总}
\label{tab:exceed_summary}
\begin{tabularx}{\textwidth}{CCCCC}
\toprule
指标 & 零高度 $P_{exceed}$ & 零高度等级 & 最大动压 $P_{exceed}$ & 最大动压等级 \\
\midrule
水平脱靶量 $D_h$ & $3.80\%$ & 低风险 & $4.80\%$ & 低风险 \\
三维脱靶量 $D_{3D}$ & $3.80\%$ & 低风险 & $0\%$ & 低风险 \\
最大角速度 $\omega_{max}$ & $0\%$ & 低风险 & $0\%$ & 低风险 \\
最高点高度 $h_{max}$ & $0\%$ & 低风险 & $0\%$ & 低风险 \\
峰值动压 $q_{max}$ & $0\%$ & 低风险 & $0\%$ & 低风险 \\
\bottomrule
\end{tabularx}
\end{table}

综合结论如下：

（1）两种工况均具备较高的命中精度。零高度逃逸和最大动压逃逸的命中风险均为低风险，超限概率分别为 $3.80\%$ 和 $4.80\%$，命中率分别为 $96.20\%$ 和 $95.20\%$，均满足 $95\%$ 置信水平的工程要求。

（2）全部非脱靶指标均为低风险。两类工况的最大角速度、最高点高度和峰值动压的超限概率均为 $0\%$，远低于各自的工程阈值，控制裕度和结构安全裕度充足。其中最大动压工况的峰值动压最大值（$49.8\ \mathrm{kPa}$）已接近 $50\ \mathrm{kPa}$ 的承载阈值，在极端扰动组合下需关注结构安全边界。

（3）最敏感风险因素差异显著。零高度逃逸主要受风场环境影响，最大动压逃逸主要受气动参数不确定性影响。这一差异由两者的飞行力学环境决定——低动压工况风扰主导，高动压工况气动偏差主导。相应地，工程上应针对零高度逃逸重点加强风场测量与预报能力，而针对最大动压逃逸重点提高气动数据的精度与置信度。

（4）最大动压逃逸的尾部风险值得关注。虽然命中率达标，但散布量级大、标准差大，最恶劣组合下的脱靶量可达 $1414.684\ \mathrm m$，显著高于零高度逃逸的 $69.820\ \mathrm m$。若需进一步提升安全裕度，可考虑缩小安全回收区的覆盖范围或增加控制系统的鲁棒性设计。
